\section{A uniform support bound for extendible orthogonal characters}
\label{sec:orthogonal-extendible-support}

\paragraph{Author check and review status.}
This is a new author derivation.  The complete written proof and
all invoked source passages have passed the author check.  The
exact author supplement also passed; it is separate from the
universal proof. Whole-proof reviews by /root and /root/g34\_eightfolds
pass, with independent exact arithmetic recorded separately. Actual
check scopes, hashes and times are recorded separately; no finite
experiment substitutes for an all-dimension argument below.

Let $V$ be a quadratic space of natural dimension $d\ge22$ over
$F=\mathbf F_q$, with nondegenerate polar form
$B(x,y)=Q(x+y)-Q(x)-Q(y)$.  In characteristic two $d$ is even.
Put $O=O(V)$ and $S=\Omega(V)=[O,O]$.  Let $\rho$ be a nontrivial
irreducible complex representation of $S$, of degree $a$, which
extends to an irreducible representation $R$ of $O$.  For a finite
projective-order complex operator $T$, define
\[
 \nu(T)=a-\max_{\lambda}\dim\ker(T-\lambda I).
 \tag{OS.1}
\]
Rescaling $T$ does not change this number.  A compatible normalizer
means an operator inducing an automorphism of $\rho(S)$.

\subsection*{D-OL-06. Two opposite abelian radicals and their weights}

Choose a hyperbolic pair $e,f$ and write
\[
 V=Fe\oplus W\oplus Ff,\qquad
 Q(xe+w+yf)=xy+Q_W(w),\qquad N=\dim W=d-2\ge20.
 \tag{OS.2}
\]
For $v\in W$ define
\[
 u_e(v)e=e,\quad
 u_e(v)w=w-B(v,w)e,\quad
 u_e(v)f=f+v-Q(v)e.
 \tag{OS.3}
\]
Direct substitution in (OS.2) proves that these are isometries and
that $u_e(v)u_e(v')=u_e(v+v')$.  Thus
$U_e=\{u_e(v):v\in W\}$ is the additive group of $W$.
Let $U_f$ be the opposite group obtained by interchanging $e$ and $f$.
Every nonidentity $u_e(v)$ has natural moving rank two.  In odd
characteristic it has odd prime order, and therefore belongs to $S$,
since $O/S$ is a $2$-group.  In characteristic two it belongs to $S$
by the Dickson-invariant criterion
\cite[Lemma~2.1, p.~383]{OS:TaylorVinroot2018}.  Consequently
$U_e,U_f\subseteq S$ in both characteristics.

We give a complete generation argument.  Put $K=\langle U_e,U_f\rangle$.
The stabilizer $P$ of $Fe$ in $O$ is $U_e\rtimes L$, where $L$ acts
by $e\mapsto ce$, $f\mapsto c^{-1}f$, and an arbitrary isometry on
$W$.  This description follows by first correcting the image of $f$
with a unique element of $U_e$ and then restricting to $W$.
The group $L$ normalizes both radicals, so $P$ normalizes $K$.
The isometry $w_0$ interchanging $e,f$ and fixing $W$ also normalizes
$K$.

For every nonsingular $z$, write
$r_z(x)=x-B(x,z)Q(z)^{-1}z$ for its reflection (or symmetry in
characteristic two).  If $z=xe+w+yf$ with $y=0$, then $r_z$ fixes
$e$, so lies in $P$.  If $y\ne0$, then
\[
 u_e(-w/y)z=(Q(z)/y)e+yf.
 \tag{OS.4}
\]
Reflection in the latter vector interchanges $Fe$ and $Ff$ with
nonzero scalar factors, so belongs to $Lw_0$.  Thus every reflection
normalizes $K$.  The full orthogonal group is generated by its
reflections, with the sole exception $O_4^+(2)$, which is outside
our range; this is \cite[Corollary~11.42, p.~159]{OS:Taylor1992},
including Theorems~11.39 and 11.41 and their characteristic-two
qualification.  Hence $K\triangleleft O$.  It is a noncentral
subgroup of $S$.  The perfectness of $S$ and simplicity of $S/Z(S)$,
\cite[Theorems~11.47--11.48, pp.~161--162]{OS:Taylor1992}, imply
$KZ(S)=S$, and then
$S=[S,S]=[KZ(S),KZ(S)]\subseteq K$.  Therefore
\[
                         \langle U_e,U_f\rangle=S.
 \tag{OS.5}
\]

The two radicals are conjugate inside $S$.  To see this without
assuming that $w_0\in S$, embed $O(W)$ in the pointwise stabilizer
of $e,f$.  It surjects onto $O/S$: in odd characteristic a binary
nondegenerate subspace represents both nonzero square classes and
its reflections generate both determinant/spinor classes; in even
characteristic one symmetry supplies the nontrivial Dickson class.
This is exactly the binary-complement calculation in D-OL-04.
Choose $l\in O(W)$ with $w_0l\in S$.  It conjugates $U_e$ to $U_f$.
Since $\rho$ is nontrivial irreducible, (OS.5) gives
\[
 b:=\dim \rho^{U_e}=\dim\rho^{U_f}\le a/2.
 \tag{OS.6}
\]
Indeed the two fixed spaces have zero intersection.

Fix a nontrivial additive character $\psi$ of $\mathbf F_p$.
The characters of $U_e$ are indexed by $z\in W$ through
\[
 u_e(v)\longmapsto
 \psi\bigl(\operatorname{Tr}_{F/\mathbf F_p} B(z,v)\bigr).
 \tag{OS.7}
\]
Their multiplicities in $\rho$ are constant on the orbits of the
embedded $\Omega(W)$ in the Levi subgroup.  These nonzero orbits
are precisely
\[
                  \mathcal O_t=\{z\in W\setminus\{0\}:Q(z)=t\},
                  \qquad t\in F.
 \tag{OS.8}
\]
Witt's extension theorem gives the assertion first for $O(W)$
\cite[Theorem~7.4, pp.~57--58]{OS:Taylor1992}.  The pointwise
stabilizer of a vector contains $O(A)$ for a nondegenerate
complement $A$ of dimension at least $N-2\ge2$ in its perpendicular
space.  The same binary-complement calculation makes this
stabilizer surject onto $O(W)/\Omega(W)$.  The $O(W)$-orbit and
$\Omega(W)$-orbit consequently coincide.  This proves (OS.8),
also for nonzero singular vectors.

Here are the needed orbit estimates with their full error bounds.
For any quadratic form $Q_0$ on a vector space $X$ of dimension $m$,
whose polar form has rank $s$, and any nontrivial additive character
$\eta$ of $F$, orthogonality gives
\[
 \left|\sum_{x\in X}\eta(Q_0(x))\right|^2
 =q^m\sum_{z\in\operatorname{rad}B_0}\eta(Q_0(z))
 \le q^{2m-s}.
 \tag{OS.9}
\]
The equality follows by setting the difference of the two summation
variables equal to $z$; the inner sum of $\eta(B_0(x,z))$ is zero
unless $z$ is in the radical.  Taking absolute values in the last
sum proves the inequality, including characteristic two and a
nonzero quadratic form on the polar radical.  Fourier inversion
therefore yields
\[
 \left|\#\{x\in X:Q_0(x)=t\}-q^{m-1}\right|
 \le q^{m-s/2}.
 \tag{OS.10}
\]
In particular, every orbit in (OS.8) has size at least
\[
                  L_N=q^{N-1}-q^{\lceil N/2\rceil}-1.
 \tag{OS.11}
\]
If a linear isometry $g$ of $W$ has moving rank two, then
$\dim\ker(g-1)=N-2$.  The polar form restricted to this fixed space
has radical of dimension at most two and hence rank at least $N-4$.
Applying (OS.10) gives
\[
 \frac{|\mathcal O_t\cap\ker(g-1)|}{|\mathcal O_t|}
 \le F_N(q):=
 \frac{q^{N-3}+q^{\lceil N/2\rceil}}
      {q^{N-1}-q^{\lceil N/2\rceil}-1}.
 \tag{OS.12}
\]
For $N\ge20$, division by $q^{N-1}$ gives the uniform bounds
\[
 F_N(q)\le\frac{q^{-2}+q^{-9}}{1-q^{-9}-q^{-19}}
 <\begin{cases}1/8,&q\ge3,\\1/3,&q\ge2.\end{cases}
 \tag{OS.13}
\]
For completeness the two inequalities follow by substituting
$q=3$ and $q=2$, respectively: the numerator decreases and the
positive denominator increases with $q$, and cross multiplication
uses only positive integers.  This is a universal inequality in
$N,q$, not a finite-range extrapolation.

Suppose now that an operator $A$ on the representation space
normalizes $U_e$ and induces an action of prime order $l$ on its
weight set.  For each length-$l$ weight orbit the corresponding
weight spaces have equal dimension $m$; an eigenspace of $A$ on
their direct sum has dimension at most $m$.  This follows by
successively recovering its $l-1$ other components from the first.
No assumption that $A^l$ is scalar is required.  If in each nonzero
orbit (OS.8) at most a proportion $f_0$ of the weights are fixed,
the constant multiplicities and (OS.6) give
\[
                \nu(A)\ge (1-1/l)(1-f_0)(a-b)
                         \ge\tfrac12(1-1/l)(1-f_0)a.
 \tag{OS.14}
\]
In particular, for a prime-order isometry of $V$ with moving rank
two which fixes a hyperbolic pair, (OS.13)--(OS.14) give
\[
 \nu(R(g))\ge7a/32\quad(q\text{ odd}),\qquad
 \nu(R(g))\ge2a/9\quad(q\text{ even},\ l\text{ odd}).
 \tag{OS.15}
\]
These estimates are valid for $g\in O$, not only for $g\in S$;
the irreducibility needed in (OS.6) concerns $R|_S$.

Every rank-two isometry $g$ used here does fix such a hyperbolic
pair.  Embed $\operatorname{im}(g-1)$ in a nondegenerate subspace
$A$ of dimension at most four: choose dual partners for the radical
of its restricted polar form.  Then $A$ is $g$-stable, and
$A^\perp\subseteq\ker(g-1)$ because
$\operatorname{im}(g-1)^\perp=\ker(g-1)$.
The nondegenerate space $A^\perp$ has dimension at least $d-4\ge18$
and contains a hyperbolic pair over a finite field
\cite[Theorem~11.2 and its ensuing decomposition, p.~138]{OS:Taylor1992}.

\subsection*{D-OL-07. Linear and semilinear normalizers}

We first specify completely why the normalizers below are covered
by natural semilinear similitudes.  For $d\ge22$, the orders of the
outer automorphism groups in
\cite[Section~2.2.2, Table~2.1, p.~14]{HES:LiXia2022} agree with
the following orders obtained from these similitudes.  Write
$q=p^f$, and use $\epsilon=1$ or $-1$ for the plus or minus type
when $d=2m$:
\[
\begin{array}{c|c}
\text{natural type}&|\Gamma\!GO(V)/(F^*\Omega(V))|\\ \hline
d\text{ odd},\ q\text{ odd}&2f\\
d\text{ even},\ q\text{ even}&2f\\
d=2m,\ q\text{ odd},\ q^m\not\equiv\epsilon\pmod4&4f\\
d=2m,\ q\text{ odd},\ q^m\equiv\epsilon\pmod4&8f.
\end{array}
 \tag{OS.16}
\]
Here the notation in the numerator means all semilinear similitudes,
and the quotient means its induced outer image; scalar intersections
are taken in computing the order.

We justify the calculation and its faithfulness.  One has
$|O:S|=4$ for odd $q$ and $2$ for even $q$.
For odd natural dimension a multiplier must be a square by taking
determinants; for even characteristic every multiplier is a square.
In these cases $GO=F^*O$.  For even dimension and odd $q$ the
multiplier map is surjective.  On hyperbolic planes scale the first
coordinate by the desired multiplier; on an anisotropic binary
summand use multiplication in $\mathbf F_{q^2}$ with the prescribed
norm.  The norm is surjective since finite-field multiplicative
groups are cyclic.  In odd characteristic and even dimension,
$-I\in S$ exactly when $q^m\equiv\epsilon\pmod4$: its determinant
is one and its spinor norm is the discriminant of $Q$, namely
$(-1)^m$ times a square in plus type and times a nonsquare in minus
type.  These facts give the linear factors $2,2,4,8$ in (OS.16).
The determinant/spinor description is
\cite[Section~2, pp.~41--43]{OL:KTV2023}; the reflection spinor
calculation is also \cite[p.~163]{OS:Taylor1992}.

Each field automorphism occurs in a semilinear isometry.  This
follows directly from the finite-field quadratic classification
in \cite[pp.~138--140]{OS:Taylor1992}: the split form has coefficients
in the prime field.  For odd characteristic the nonsplit binary
form may be written $x^2-\epsilon y^2$ with $\epsilon$ a nonsquare;
Frobenius is corrected by scaling the $y$ coordinate by
$\epsilon^{(p-1)/2}$.  In characteristic two write it as
$x^2+xy+\delta y^2$, where $\operatorname{Tr}(\delta)=1$.
Since $\delta^p-\delta$ has trace zero, choose $c$ with
$c^2+c=\delta^p-\delta$ and correct Frobenius by
$(x,y)\mapsto(x+cy,y)$.  In odd natural dimension rescale the form
so that its final one-dimensional coefficient is one.  Thus there
are all $f$ field components, even if a chosen lift has order
$2f$ rather than $f$.

The centralizer of $S$ in the semilinear group consists only of
linear scalars.  Indeed, an element centralizing $S$ preserves
every space $\operatorname{im}(u_e(v)-1)=\langle e,v\rangle$;
their intersection is $Fe$.  As $e$ is arbitrary, it fixes every
singular projective point.  With two mutually perpendicular singular
basis vectors $e_1,e_2$, fixing $F(e_1+e_2)$ first makes their
scalar factors equal, and fixing all $F(e_1+te_2)$ then forces its
field automorphism to fix every $t\in F$.  It is now linear.  Three
hyperbolic pairs and the singular vectors $e_i+f_j$ for $i\ne j$
force a common scalar on the hyperbolic part.  For a vector $z$ of
the anisotropic part, the singular vector
$e_1+z-Q(z)f_1$ forces the same scalar on $z$.
The same conclusion holds if the element acts trivially on
$S/Z(S)$: its commutators with $S$ are central and define a
homomorphism from the perfect group $S$ to $Z(S)$, hence are trivial.
Consequently the image of the semilinear similitudes in
$\operatorname{Aut}(S/Z(S))$ has precisely the outer order (OS.16).
Comparison with the published complete table proves that it is the
whole automorphism group.  Dimension eight triality does not occur.

Finally, a nontrivial irreducible $S$-representation has central
kernel.  The center of the natural $S$ has order at most two, so
every subgroup of it is preserved by the above automorphisms.
An automorphism of $\rho(S)$ and the induced natural lift which
agree on the simple quotient differ by a central homomorphism;
perfectness again makes that difference trivial.  Thus every
compatible normalizer $T$ is implemented on $S$ by a natural
semilinear similitude $H$.

There is an important extension issue.  The chosen extension $R$
need not itself be invariant under conjugation by $H$.  Instead
Schur's lemma gives a linear character $\kappa_H$ of $O/S$ such that
\[
 T R(x)T^{-1}=\kappa_H(x)R(HxH^{-1})\qquad(x\in O).
 \tag{OS.17}
\]
To prove this, both operators on the right and left induce the
same conjugation on $\rho(S)$; their ratio is scalar.  Multiplication
shows that these scalars form a character, and the character is
trivial on $S$.  For a reflection $r$, $\kappa_H(r)^2=1$.
For a linear $H=h$, it follows that
\[
 [T,R(r)]=\kappa_h(r)R([h,r]),\qquad
 \nu(R([h,r]))\le\operatorname{rank}([T,R(r)]-I)\le2\nu(T).
 \tag{OS.18}
\]
The last inequality follows by choosing $\lambda$ with
$\operatorname{rank}(T-\lambda I)=\nu(T)$ and writing
$TR(r)-R(r)T=(T-\lambda I)R(r)-R(r)(T-\lambda I)$.
The first inequality uses projective support, so the scalar
$\kappa_h(r)$ is harmless.  No equality of ordinary moving ranks
after discarding that scalar is asserted.

\paragraph{Linear normalizers in odd characteristic.}
If $h$ is not scalar, some reflection $r_v$ does not commute with
$h$, because reflections generate $O$ and its natural centralizer
is scalar.  Then
$k=[h,r_v]=r_{hv}r_v\ne1$ is an isometry of determinant one and
natural moving rank at most two.  A rank-one isometry in odd
characteristic is a reflection and has determinant $-1$, so $k$
has rank two.  The same argument applies to every nonidentity
power of $k$.  Choose a nonidentity prime-order power $g$ of $k$.
Its rank is two.  Equation (OS.15) and
$\nu(R(k))\ge\nu(R(g))$ imply from (OS.18) that
\[
                         \nu(T)\ge7a/64.
 \tag{OS.19}
\]
The rank-one assertion can also be checked directly: an isometry
$I+v\ell$ of rank one has $\ell=cB(v,-)$, and preservation of $Q$
gives $cQ(v)=-1$, so its determinant is $1+2cQ(v)=-1$.
The element $k$ need not belong to $S$; use of the actual extension
$R$ in (OS.15) is essential here.

\paragraph{Linear normalizers in characteristic two.}
Rescale $h$ to an isometry.  If some nonsingular $v$ satisfies
$B(v,hv)\ne0$, the plane $\langle v,hv\rangle$ is nondegenerate.
The product $k=r_{hv}r_v$ is a nonidentity rotation on this plane
and is the identity on its perpendicular space.  In a hyperbolic
plane these rotations have order dividing $q-1$; in an anisotropic
plane they are multiplication by norm-one elements of
$\mathbf F_{q^2}^*$, of order dividing $q+1$.  Hence a nonidentity
prime-order power $g$ has odd order and rank two.  Equations
(OS.15) and (OS.18) give
\[
                         \nu(T)\ge a/9>7a/64.
 \tag{OS.20}
\]

It remains to treat the case $B(v,hv)=0$ for every nonsingular $v$.
Put $F_0(v)=B(v,hv)$, a homogeneous quadratic polynomial.  We show
that it vanishes on every vector.  If $q\ge4$, for any nonzero
singular $x$ choose a nonsingular $y\in x^\perp$.  Such a $y$
exists because the polar rank of $x^\perp$ is $d-2>0$.
For every $t\ne0$, $x+ty$ is nonsingular, and
$0=F_0(x+ty)=F_0(x)+tB_{F_0}(x,y)$.
Two distinct nonzero $t$ give $F_0(x)=0$.
For $q=2$, the nondegenerate quadratic quotient
$x^\perp/Fx$ contains an anisotropic plane: this is immediate if
its anisotropic kernel is a plane, and in two hyperbolic planes
one may take $y=e_1+f_1$ and $z=e_1+e_2+f_2$.
Lift its vectors to $x^\perp$; then
$Q(y)=Q(z)=Q(y+z)=1$.  The three equations for
$x+y,x+z,x+y+z$ give
$B_{F_0}(x,y)=B_{F_0}(x,z)=B_{F_0}(x,y+z)=F_0(x)$;
bilinearity forces $F_0(x)=0$.

Polarizing the now identically zero $F_0$ gives $h^*=h$ for the
adjoint relative to $B$.  As $h$ is an isometry, $h^*=h^{-1}$,
so $h^2=I$.  Moreover $B(x,(h-I)y)$ is alternating, since its
diagonal values are $F_0(x)=0$.  Its rank, which equals
$\operatorname{rank}(h-I)$, is even.  The Dickson criterion puts
$h$ in $S$.  Unless $T$ is scalar, $h\ne I$ and is a noncentral
unipotent element.  The known bound
\cite[Theorem~2.4(i), with conventions on pp.~2--4]{HES:Liebeck2017}
gives $|\operatorname{tr}\rho(h)|\le3a/4$ for $q=2,4$ and at most
$a/(\sqrt8-1)<9a/16$ for $q=8$.
Here $S$ is the simply connected finite group in characteristic
two in these large dimensions, so the source hypotheses apply.
Since $T$ is a scalar times $\rho(h)$ by Schur's lemma,
\[
                         \nu(T)\ge a/8>7a/64.
 \tag{OS.21}
\]
The subtraction in $\sqrt8-1$ is outside the radical.  The
comparison with $9/16$ follows from $625<648$.

\paragraph{General semilinear normalizers for $q=4,8,9$.}
These are the only nonprime fields below sixteen.  Their field
degrees $f=2,3,2$ are prime.  If the field part of $H$ is trivial,
the linear argument applies.  Otherwise write its field
automorphism as $\sigma$, of order $f$, with fixed field $F_0$.
If $T^f$ is nonscalar, its natural implementing map $H^f$ is linear,
so (OS.19)--(OS.21) give the same bound for $T$.

Suppose instead that $T^f$ is scalar.  Its conjugation on $\rho(S)$
is trivial.  Projective faithfulness, perfectness, and the centralizer
calculation above imply $H^f=cI$ on $V$ for some $c\in F^*$.
Commuting $H$ with $H^f$ shows $\sigma(c)=c$.  Surjectivity of the
norm $F^*\to F_0^*$ lets us rescale $H$ so that $H^f=I$.
If $Q(Hv)=\mu\sigma(Q(v))$, then the norm of $\mu$ is one.
The image of $b\mapsto\sigma(b)/b$ has size
$(q-1)/(|F_0|-1)$ and equals the norm-one subgroup.  Choose $b$
with $\sigma(b)/b=\mu$ and replace $Q$ by $bQ$.
The isometry groups do not change, and now
\[
                      Q(Hv)=\sigma(Q(v)),\qquad H^f=I.
 \tag{OS.22}
\]

We include the vector descent proof; the case is not assumed to
be a pure field automorphism in the original coordinates.
For every $t\in F$ and $v\in V$, the vector
$\sum_{j=0}^{f-1}H^j(tv)$ is fixed by $H$.
If the fixed vectors did not span $V$ over $F$, a nonzero $F$-linear
functional $\lambda$ annihilating their span would give
\[
       \sum_{j=0}^{f-1}\sigma^j(t)\lambda(H^jv)=0
                  \quad\text{for all }t\in F,\ v\in V.
 \tag{OS.23}
\]
The distinct field automorphisms $1,\sigma,\ldots,\sigma^{f-1}$
are linearly independent as functions to $F$: evaluate a linear
relation on $1,z,\ldots,z^{f-1}$ for an element $z$ with distinct
conjugates, and use the Vandermonde determinant.  Equation (OS.23)
would therefore force $\lambda(v)=0$ for every $v$, a contradiction.
Choose an $F$-basis of fixed vectors.  The fixed vectors are then
exactly its $F_0$-linear span, and (OS.22) says that the quadratic
form on this $d$-dimensional $F_0$-space is nondegenerate.
Choose a hyperbolic pair over $F_0$ and apply (OS.2) in this basis.
The induced action on $U_e$ parameters, and on their dual weight
vectors, is $\sigma$ entrywise.

There are at most $|F_0|^N$ fixed weight vectors in all of $W$.
Consequently each orbit (OS.8) has fixed proportion at most
\[
 \frac{|F_0|^N}{L_N}
 \le\frac{q\,|F_0|^{-N(f-1)}}{1-q^{-9}-q^{-19}}<1/4
       \quad(N\ge20,\ q\in\{4,8,9\}).
 \tag{OS.24}
\]
The inequality follows already at $N=20$ in each of the three
displayed fields, and the numerator decreases with $N$ while
the exact denominator bound is fixed.  It is an exact rational
comparison.  Norm levels not fixed by $\sigma$ have no fixed
weight vectors at all.  Applying (OS.14) with $l=f$ gives
\[
                         \nu(T)\ge3a/16>7a/64.
 \tag{OS.25}
\]
This proves the desired support bound for all semilinear
normalizers, including diagonal--field and graph--field mixtures,
not merely for pure field lifts.

\begin{theorem}[D-OL-08: closing the ordinary orthogonal dimension escape]
\label{D-OL-08}
For $d\ge22$ and $q\in\{2,3,4,5,7,8,9,11,13\}$, every
nonscalar compatible normalizer of a nontrivial irreducible
$\Omega(V)$-representation which extends irreducibly to $O(V)$
satisfies $\nu(T)\ge7a/64$.

In the actual rank-nine six-dimensional multiplicity variation
setting of D-OL-02--05, suppose the inducing image has a normal
irreducibly acting central product of quasisimple orthogonal
factors, permuted transitively.  Then their natural dimension obeys
\[
                             d\le4\cdot10^{30}-3.
 \tag{OS.26}
\]
This conclusion allows all compatible scalars, outer automorphisms
and factor permutations, and allows the inducing character field
to be larger than the original imaginary quadratic field.  It is
a bounded-dimension reduction, not a classification or an existence
assertion for the surviving parameters.
\end{theorem}

\begin{proof}
The first assertion is (OS.19)--(OS.25), since the remaining fields
in the displayed set are prime.  To prove (OS.26), suppose instead
that $d\ge4\cdot10^{30}-2$.  D-OL-05 forces all factor characters
to be ordinary, to extend irreducibly to their full orthogonal
groups, to be self-dual, and to have $q$ in the displayed set.
Their common degree is $a$ and the degree of the inducing
representation is $e=a^t$.

For an element fixing all factors, the induced operator is, up to
scalar, a tensor product of compatible factor intertwiners.  A
nonscalar factor has support at least $7/64$ of its dimension;
eigenvalue multiplicities in a tensor product are convolutions,
so the same fraction bounds the support of the whole operator.
If an element moves factors in a cycle of length $l\ge2$, conjugate
the corresponding weighted cyclic permutation to one with a
single twist.  This twist has finite projective order and an
eigenbasis.  On tensor words in that basis there are at most $a$
constant words with orbits of length one; every other word orbit
has length at least two, and on an orbit of length $b$ an
eigenspace has dimension at most one.  Thus the largest eigenspace
on that cycle has dimension at most
\[
                    a+(a^l-a)/2\le3a^l/4.
 \tag{OS.27}
\]
Tensoring the remaining cycles preserves this upper fraction.
Hence every nonscalar inertia in the inducing image has projective
support at least $7e/64$.  A nonidentity scalar inertia has actual
moving rank $e$, so satisfies the required moving-rank bound too.

The normal-core invariant bilinear form makes the inducing image
essentially self-dual.  Its multiplier is nontrivial: if it were
trivial the inducing character would be real-valued, whereas its
character field contains the original imaginary quadratic field
by the induction formula.  This is the same actual form and field
argument as D-OL-05; no equality of those fields is assumed.
The actual Hurwitz factorization has genus zero, at least twelve
genuine branch points, and total moving rank $2e+6$.
At least two inertias have nontrivial multiplier because the
inertias generate and have product one.  Each then has moving
rank at least $e/2$, by pairing the $1$- and multiplier-eigenspaces.
The remaining at least ten have moving rank at least $7e/64$.
Therefore
\[
                 2e+6\ge e+10(7e/64),\qquad e\le64.
 \tag{OS.28}
\]
But the $k=1$ case of the published degree estimate
\cite[Theorem~3.7, pp.~165--167]{OL:GLT2024} gives
$a\ge q^{d-4}(q-1)\ge2^{18}>64$ for our nontrivial ordinary
factor and $d\ge22$.  Since $e\ge a$, this is a contradiction.
This proof uses the rank-nine assumption only through the already
proved actual intermediate Hurwitz factorization.  It never
infers dominance from a source dimension inequality.
\end{proof}
